\section{Test rigs, bearings and kinematics}
\label{sec:data}

\begin{table}[htbp]
\centering
\caption{Test rigs and recordings used in this study. Bearings are physical bearings, not fault classes.}
\label{tab:datasets}
\small
\setlength{\tabcolsep}{4pt}
\begin{tabular}{lllrrl}
\toprule
Dataset & Bearings used & $f_s$ (kHz) & Record (s) & Conditions & Role \\
\midrule
Paderborn (PU) \cite{lessmeier2016} & 6 healthy, 5 outer, 6 inner (real) & 64 & 4 & 4 (3 adapt.) & adaptation, gating, slip \\
HUST bearing \cite{thuan2023hust} & 5 types $\times$ \{healthy, inner, outer\} & 51.2 & 10 & 3 loads & second-rig replication \\
CWRU \cite{smith2015} & 6205 drive end, 6203 fan end & 12 / 48$^{a}$ & $\approx$10 & 4 loads & slip, contamination \\
UORED-VAFCLS \cite{sehri2023uored} & 20 bearings, natural faults & 42 & 10 & 1 & gate safety (suppl.) \\
JNU \cite{li2013jnu} & as published & 50 & 20 & 3 speeds & reproduction only \\
\bottomrule
\end{tabular}
\begin{minipage}{\linewidth}\vspace{2pt}\footnotesize $^{a}$ Fault recordings 12~kHz; the four normal baselines used for off-rig calibration 48~kHz. Paderborn: three conditions at 1500~rpm (25~Hz) for adaptation, all four for the gate and slip analyses. JNU: the three nominal speeds of the distributed data, as used by SDALR.\end{minipage}
\end{table}


Table~\ref{tab:datasets} lists the recordings used. \textbf{Paderborn} \cite{lessmeier2016} carries the adaptation
experiments: its 6 healthy bearings and the 5 outer-race and 6 inner-race bearings of its real-damage set, each recorded 20
times at four operating conditions, are the largest public combination we know of with several physical bearings per class
\emph{and} several operating conditions per bearing, which the paired design of Section~\ref{sec:protocol} needs. We use the
three conditions that share a nominal speed of \SI{1500}{rpm} (\SI{25}{Hz}) and differ in load and radial force, denoted A1
(N15\_M01\_F10), A2 (N15\_M07\_F04) and A3 (N15\_M07\_F10); the fourth condition (N09, \SI{900}{rpm}, \SI{15}{Hz}) enters
only the gate analyses.
\textbf{HUST bearing} \cite{thuan2023hust} provides the second-rig replication: one physical bearing per class for each of
five bearing types (6204--6208) at three motor loads, giving the same three-domain, six-task structure with the same label
space. \textbf{CWRU} \cite{smith2015} and \textbf{UORED-VAFCLS} \cite{sehri2023uored} appear only in the off-rig calibration of the
gate operating point and in the supplementary slip measurement. JNU \cite{li2013jnu} is used only to reproduce the published SDALR result.

Throughout, a \emph{bearing} means a physical specimen, identified by its dataset code (for example Paderborn KA04), not a
fault class or a damage severity. The label space is $\{\text{normal},\text{inner race},\text{outer race}\}$: Paderborn's
real-damage set contains no rolling-element damage, so a ball class cannot be formed from it, and HUST's ball recordings are
excluded for comparability.

\subsection{Kinematics are resolved per bearing}
\label{sec:kinematics}

Characteristic orders follow from the ball count $n$, ball diameter $d$, pitch diameter $D$ and contact angle $\theta$
through $r=(d/D)\cos\theta$, giving $\mathrm{BPFO}=\tfrac{n}{2}(1-r)$, $\mathrm{BPFI}=\tfrac{n}{2}(1+r)$ and
$\mathrm{FTF}=\tfrac{1}{2}(1-r)$ in shaft orders. A bearing designation does not determine these values. Paderborn's
per-bearing documentation, distributed with the data as one ``Profile of rolling bearing damage'' datasheet per specimen
rather than in the dataset paper, carries the geometry on page~1 of each file. Reading all 32 datasheets gives a clean
split by manufacturer: 26 specimens state ``Pitch circle diameter mm 29.05'' (IBU and MTK) and 6 state
\SI{28.55}{\milli\metre} (FAG), all with ``Number of rolling elements pc. 8'' and rolling elements of
\SI{6.75}{\milli\metre}. The 29.05 group contains every healthy specimen K001--K006 and the real-damage bearings KA16,
KA22, KA30, KI04, KI14, KI17 and KI18 (e.g.\ \texttt{data/pu/K001/K001.pdf}); the 28.55 group contains KA04, KA15, KI16
and KI21 (e.g.\ \texttt{data/pu/KA04/KA04.pdf}) together with KB24 and KB27, which this study excludes. Per-specimen
geometry and damage mode are listed in Supplementary~S3; the damage modes are used in Section~\ref{sec:limitations}. The corresponding
outer-race orders are 3.0706 and 3.0543, neither of which equals the 3.0531 of the CWRU SKF 6203 that is frequently reused
for Paderborn. We compute kinematics separately for each specimen, and no gate is computed for a bearing whose geometry its own rig does
not publish. HUST is such a case: its dataset paper gives bore, outside and ball diameters and the ball count, but not the
pitch diameter and no characteristic orders \cite{thuan2023hust}. Substituting the nominal mean diameter $(d_i+d_o)/2$
would supply the missing number, and it is worth seeing what that costs: on Paderborn, where the true value is documented,
the same substitution gives \SI{28.50}{\milli\metre} against \SI{29.05}{\milli\metre} and shifts the outer-race order by
\SI{0.58}{\percent}, a systematic error larger than the 5--95\,\% range of slip we measure on that geometry
(\SI{0.53}{\percent}, Supplementary~S5). The gate's $\pm\SI{2}{\percent}$ slip window would absorb such an error, but only by
spending its tolerance on a geometry guess rather than on slip, and a guessed geometry would make every HUST verdict
untraceable. We therefore use HUST for raw windows only and report no HUST gating result. Shaft speed is taken per
recording from the measured speed channel on Paderborn and from the in-file value on CWRU.

\subsection{What a window can represent}
\label{sec:windows}

The adaptation experiments use SDALR's input contract: windows of 2048 samples (Section~\ref{sec:protocol}). At
\SI{64}{kHz} and \SI{25}{Hz} shaft speed a window lasts \SI{32}{ms}, 0.8 shaft revolutions and about 2.5 outer-race
impacts; its envelope spectrum has a resolution of \SI{31}{Hz}, so the first outer-race (\SI{77}{Hz}) and inner-race
(\SI{123}{Hz}) lines lie about 1.5 bins apart. A window can carry the impulsiveness and energy of a fault but not its
kinematic signature. The gates, by contrast, use the first 249\,600 samples of each recording (\SI{3.9}{s}, resolution
\SI{0.26}{Hz}, about twelve bins inside the slip window of the first outer-race harmonic). Section~\ref{sec:res-mechanism}
tests what the difference is worth with a classifier built on the gates' representation.
